\section{Codes and the module extension}
\label{sec:module}

Fix the labelled graphs and the group \(\Gamma\) supplied by
Proposition~\ref{prop:data}.  Retain the notation
\[
  n_i=|\operatorname{Vert}(\Theta_i)|,\qquad
  k_i=\lfloor n_i/2\rfloor,\qquad \eta=\frac1{100}.
\]
In particular, \(n_i>100\).  A copy \(P=P_{g,i}\) carries its prescribed
identification with \(\Theta_i\), and its vertex set is a subset of \(\Gamma\).
For every choice of root and copy to which Lemma~\ref{lem:overlap} applies,
the ball furnished by that lemma is an induced tree with fewer than
\(\eta n_i\) vertices, by Proposition~\ref{prop:data}.

\subsection{Port patterns and codes}

We call an outgoing generator label a \emph{port}.
At a vertex \(v\) of \(\Theta_i\), let
\[
  J_i(v)\subseteq S^{\pm1}
\]
be the set of outgoing labels of its incident edges.  Reducedness of the
labelling and \(d\)-regularity give \(|J_i(v)|=d\).  Choose a prime \(p>2\)
large enough that the \(d\)-element subsets of \(S^{\pm1}\) admit an
injection
\[
  J\longmapsto a_J\in\F_p^\times.
\]
Define the decoder
\[
  \mathcal J(a)=
  \begin{cases}
    J,&a=a_J\text{ for a \(d\)-element subset }J\subseteq S^{\pm1},\\
    \varnothing,&\text{otherwise},
  \end{cases}
  \qquad a\in\F_p.
\]
The prototype pattern on \(\Theta_i\) is
\[
  w_i(v)=a_{J_i(v)}.
\]
Every coordinate of \(w_i\) is nonzero.  Write \(J_P(h)\) and \(w_P(h)\)
for the corresponding data transported to a copy \(P=P_{g,i}\), so that
\(\mathcal J(w_P(h))=J_P(h)\).

On a finite coordinate space \(\F_p^E\), orthogonal complements will
always refer to the nondegenerate bilinear form
\[
  \langle z,z'\rangle=\sum_{v\in E}z(v)z'(v)\in\F_p.
\]
For \(z\in\F_p^E\), let \(\supp z=\{v\in E:z(v)\ne0\}\).

We seek a linear code containing the prescribed word $w_i$ with two
different support guarantees. Any other codeword must differ from $w_i$
at many vertices, forcing many missing prescribed edge incidences after
decoding. Vectors in the dual code will instead supply the module
relations. Their support must be larger than an overlap ball, so that a
nonzero relation on a farthest copy cannot cancel against all the other
relations. The following Gilbert--Varshamov type counting argument
\cite{Gilbert1952,Varshamov1957} obtains both bounds while keeping $w_i$
in the code.

\begin{lemma}[Codes with both distances]
\label{lem:codes}
For every \(i\) there is a \(k_i\)-dimensional subspace
\(C_i\subseteq\F_p^{\operatorname{Vert}(\Theta_i)}\) such that \(w_i\in C_i\)
and
\[
  |\supp z|>\eta n_i
  \qquad
  \text{for every nonzero }z\in C_i\cup C_i^\perp.
\]
\end{lemma}

\begin{proof}
Write \(n=n_i\), \(k=k_i\), and \(w=w_i\), and choose \(C\) uniformly among
the \(k\)-dimensional subspaces containing \(\langle w\rangle\).
For a fixed \(z\notin\langle w\rangle\), the quotient by
\(\langle w\rangle\) gives
\begin{equation}
\label{eq:primal-probability}
  \Pr(z\in C)
  =\frac{p^{k-1}-1}{p^{n-1}-1}
  <p^{k-n}.
\end{equation}
Indeed, a uniform \(a\)-dimensional subspace of an \(N\)-dimensional
vector space contains a fixed nonzero vector with probability
\((p^a-1)/(p^N-1)\), by counting nonzero vectors and symmetry.
Nonzero vectors of \(\langle w\rangle\) have full support.

Orthogonal complementation is a bijection from the sample space just
used to the \((n-k)\)-dimensional subspaces of the \((n-1)\)-dimensional
space \(w^\perp\).  Consequently, for a fixed nonzero \(z\),
\begin{equation}
\label{eq:dual-probability}
  \Pr(z\in C^\perp)=
  \begin{cases}
    \displaystyle\frac{p^{n-k}-1}{p^{n-1}-1}<p^{1-k},
       &z\in w^\perp,\\[6pt]
    0,&z\notin w^\perp.
  \end{cases}
\end{equation}
This argument uses nondegeneracy of the form on the ambient coordinate
space, not of its restriction to \(w^\perp\); it remains valid when
\(\langle w,w\rangle=0\).

We bound the number of low-weight vectors.  Put \(r=\lfloor\eta n\rfloor\)
and \(q=\eta/(1-\eta)<1\).  The binomial theorem gives
\[
  q^r\sum_{j=0}^r\binom nj
  \le \sum_{j=0}^r\binom nj q^j
  \le (1+q)^n.
\]
It follows that
\[
  \sum_{j=0}^r\binom nj
  \le q^{-\eta n}(1+q)^n
  =2^{H_2(\eta)n},
  \qquad
  H_2(t)=-t\log_2t-(1-t)\log_2(1-t).
\]
Here \(H_2(1/100)<9/100\): use \(\log_2 100<7\),
\(-\log(1-\eta)\le\eta/(1-\eta)\), and \(\log2>1/2\).
Since \(p\ge2\), the Hamming ball therefore satisfies
\begin{equation}
\label{eq:hamming-ball}
  \bigl|\{z:|\supp z|\le r\}\bigr|
  =\sum_{j=0}^r\binom nj(p-1)^j
  \le p^{\eta n}2^{H_2(\eta)n}
  \le p^{(\eta+H_2(\eta))n}
  <p^{n/4}.
\end{equation}
Every nonzero vector in this ball lies outside \(\langle w\rangle\).
By \eqref{eq:primal-probability}--\eqref{eq:hamming-ball} and the union
bound, the probability that either \(C\) or \(C^\perp\) contains such a
vector is less than
\[
  p^{n/4}\bigl(p^{k-n}+p^{1-k}\bigr)
  \le p^{-n/4}+p^{3/2-n/4}<1,
\]
where \(k=\lfloor n/2\rfloor\) and \(n>100\).
Thus a suitable \(C\) exists.
\end{proof}

Choose a code \(C_i\) as in Lemma~\ref{lem:codes} for every \(i\).
Transporting coordinates defines subspaces \(C_P\) and \(C_P^\perp\)
on every copy \(P=P_{g,i}\).  The transport preserves the bilinear
form, so these subspaces are orthogonal complements in \(\F_p^P\).

\subsection{The quotient module}

Translation-invariant quotients of finite-field lamp modules also occur
in Austin's irrational-kernel construction
\cite[Sections~1--2]{Austin2013}. Here the relations are dual codewords on
all translated expander copies. We must show that imposing them globally
creates exactly the prescribed relations on each individual copy.

Let \(\F_p^{(\Gamma)}\) denote the finite-support coordinate space with
basis \((\mathbf e_h)_{h\in\Gamma}\), on which \(\Gamma\) acts by
\(g\mathbf e_h=\mathbf e_{gh}\).  Regard \(\F_p^P\) as its coordinate
subspace by extension by zero.  Define the algebraic sum
\begin{equation}
\label{eq:relation-module}
  W_0=\sum_{\substack{g\in\Gamma\\i\ge1}} C_{P_{g,i}}^\perp
  \subseteq\F_p^{(\Gamma)}.
\end{equation}
Every vector in this sum is a sum of finitely many terms.  Left
translation carries \(C_{P_{g,i}}^\perp\) to
\(C_{P_{hg,i}}^\perp\), so \(W_0\) is \(\Gamma\)-invariant.

\begin{lemma}[Exact coordinate relations]
\label{lem:module-intersection}
For every copy \(P\),
\begin{equation}
\label{eq:module-intersection}
  W_0\cap\F_p^P=C_P^\perp.
\end{equation}
\end{lemma}

\begin{proof}
The inclusion from right to left follows from
\eqref{eq:relation-module}.  For the converse, let
\[
  z=\sum_{Q\in\mathcal F}z_Q\in W_0\cap\F_p^P,
  \qquad 0\ne z_Q\in C_Q^\perp,
\]
where \(\mathcal F\) is finite and its copies are distinct.  Such an
expression is obtained by combining terms on the same copy and
discarding zero terms.

Suppose that \(\mathcal F\) contains a copy different from \(P\).
Root the coned graph at \(R=c_P\), and choose
\(Q\in\mathcal F\setminus\{P\}\) whose center is farthest from \(R\).
Every other copy in \(\mathcal F\), and also \(P\) itself, has center at
distance from \(R\) at most \(d_{\widehat X}(c_Q,R)\).
Lemma~\ref{lem:overlap} places all their vertex intersections with \(Q\)
in one intrinsic ball \(B_Q\subseteq Q\).
If \(Q\) has type \(j\), Proposition~\ref{prop:data} gives
\[
  |B_Q|<\eta n_j<|\supp z_Q|.
\]
Choose \(h\in\supp z_Q\setminus B_Q\).  Then \(h\notin P\) and
\(h\) belongs to no other copy in \(\mathcal F\).  The \(h\)-coordinate
of \(z\) is therefore \(z_Q(h)\ne0\), contradicting
\(z\in\F_p^P\).
Hence the expression has no term away from \(P\), and
\(z\in C_P^\perp\).
\end{proof}

\begin{definition}
Define the discrete additive group and distinguished element
\[
  V=\F_p^{(\Gamma)}/W_0,\qquad v_*=[\mathbf e_1]\in V,
\]
with the induced \(\Gamma\)-action, and put
\[
  \Gsur=V\rtimes\Gamma.
\]
We use the multiplication
\((v,g)(v',g')=(v+gv',gg')\) in \(\Gsur\).
\end{definition}

For \(h\in\Gamma\), one has \(hv_*=[\mathbf e_h]\).
The group \(\Gsur\) is countable: \(\Gamma\) is finitely generated, and the
finite-support coordinate space over the finite field \(\F_p\), its
quotient \(V\), and \(V\times\Gamma\) are countable.
In fact \(\Gsur\) is finitely generated.  The additive group of \(V\) is
generated by the \(\Gamma\)-orbit of \(v_*\), since these are the images
of the coordinate basis and multiplication by a scalar in \(\F_p\)
is repeated addition.  In the semidirect product,
\[
  (0,h)(v_*,1)(0,h)^{-1}=(hv_*,1).
\]
It follows that \((v_*,1)\) together with the finitely many lifts
\((0,s)\), \(s\in S\), generates \(\Gsur\).

\subsection{The compact dual and its finite marginals}

Write the compact dual of \(V\) as
\[
  \Omega=\operatorname{Hom}_{\F_p}(V,\F_p).
\]
The corresponding complex character is
\(v\mapsto \exp(2\pi\mathrm i\,\omega(v)/p)\), where field elements
are identified with residues modulo \(p\).
We equip \(\Omega\) with the topology of pointwise convergence,
equivalently the subspace topology from \(\F_p^V\).
It is a closed subgroup of this compact product, and hence a compact
abelian group.  Let \(m_\Omega\) be its Haar probability measure.
For \(\omega\in\Omega\) and \(h\in\Gamma\), define
\[
  X_\omega(h)=\omega(hv_*).
\]
The notation \(X_\omega\) denotes a coordinate pattern, not the
Cayley graph \(X\).

\begin{lemma}[Uniform pattern marginals]
\label{lem:marginals}
For a copy \(P=P_{g,i}\), the restriction map
\[
  \operatorname{res}_P:\Omega\longrightarrow\F_p^P,
  \qquad \omega\longmapsto X_\omega|_P
\]
has image exactly \(C_P\).  Its pushforward of \(m_\Omega\) is the
uniform measure on \(C_P\).  In particular, for every \(a\in C_P\),
\begin{equation}
\label{eq:uniform-marginal}
  m_\Omega\{\omega:X_\omega|_P=a\}=p^{-k_i},
\end{equation}
and this applies to \(a=w_P\).
\end{lemma}

\begin{proof}
Let \(q_P:\F_p^P\to V\) be the restriction of the quotient map.
Lemma~\ref{lem:module-intersection} gives
\[
  \ker q_P=C_P^\perp.
\]
For \(z\in\F_p^P\) and \(\omega\in\Omega\),
\[
  \omega(q_P(z))
  =\sum_{h\in P}z(h)\omega(hv_*)
  =\langle z,X_\omega|_P\rangle.
\]
Thus \(X_\omega|_P\) annihilates \(C_P^\perp\), and belongs to
\((C_P^\perp)^\perp=C_P\).

Conversely, if \(a\in C_P\), the rule
\[
  q_P(z)\longmapsto\langle z,a\rangle
\]
is a well-defined linear functional on \(q_P(\F_p^P)\).
Extend a basis of this subspace to a basis of \(V\), and extend the
functional linearly to obtain \(\omega\in\Omega\).
Then \(X_\omega|_P=a\), proving surjectivity.
No continuity condition is lost in this extension, since \(V\) is
discrete.

The map \(\operatorname{res}_P\) is a continuous homomorphism of compact
groups onto the finite group \(C_P\).  Its pushforward of Haar
probability is translation-invariant and hence uniform on \(C_P\).
As \(|C_P|=p^{k_i}\), formula~\eqref{eq:uniform-marginal} follows.
\end{proof}

\begin{remark}
Lemma~\ref{lem:marginals} specifies each individual copy marginal.
It makes no independence assertion about patterns on different
copies.  The exact intersection
\eqref{eq:module-intersection}, which uses simultaneous containment
of all relevant overlaps in one ball, is what guarantees these
marginals despite the relations imposed on every translated copy.
\end{remark}
